\documentclass[10pt,twocolumn]{article}

% ------- arXiv-standard packages -------
\usepackage[T1]{fontenc}
\usepackage[utf8]{inputenc}
\usepackage{lmodern}
\usepackage[letterpaper, margin=1in, columnsep=0.25in]{geometry}
\usepackage{amsmath,amssymb,amsthm}
\usepackage{graphicx}
\usepackage{booktabs}
\usepackage{array}
\usepackage{multirow}
\usepackage{xcolor}
\usepackage{url}
\usepackage{hyperref}
\usepackage{microtype}
\usepackage{listings}
\usepackage{algorithm}
\usepackage{algpseudocode}
\usepackage{caption}
\usepackage{subcaption}
\usepackage{cite}
\usepackage{balance}

\hypersetup{
  colorlinks=true,
  linkcolor=blue!70!black,
  citecolor=green!50!black,
  urlcolor=blue!70!black,
}

\lstset{
  basicstyle=\ttfamily\small,
  breaklines=true,
  columns=fullflexible,
}

% --------------------------------------------------
\title{\textbf{Breaking the Byte Barrier: Why 16-Bit Native Entropy Coding\\Matters for Medical Image Compression}}

\author{
  Kuldeep Singh\\
  Innovaccer Inc.\\
  \texttt{kuldeep.singh@innovaccer.com}\\
  \url{https://github.com/pappuks/medical-image-codec}
}

\date{}
% --------------------------------------------------

\begin{document}
\maketitle

\begin{abstract}
The dominant entropy coders in modern compression systems---Huffman,
arithmetic coding, LZ77, and FSE/ANS as used in Zstandard---were designed for
8-bit byte streams. Medical imaging breaks this assumption: DICOM stores pixels
at 10--16 bits per sample, yielding symbol alphabets of 1,024 to 65,536
entries. We show that this mismatch is not merely inconvenient but structurally
costly: byte-splitting a 16-bit residual distribution discards mutual
information between the high and low bytes, inflating compressed output by
10--22\% compared to a 16-bit-native entropy coder.

We present \textbf{MIC} (Medical Image Codec), a lossless codec that addresses
this gap with four contributions:
\begin{enumerate}
  \item \textbf{16-bit-native entropy coding}: An extended FSE implementation
  supporting up to 65,535 symbols, paired with a 16-bit run-length encoder that
  captures word-level structure invisible to byte-oriented compressors. MIC
  outperforms Delta+Zstandard by 10--22\% on all tested modalities.

  \item \textbf{Multi-state ANS decoding as an ILP design axis}: A four-state
  interleaved ANS decoder that breaks the serial dependency chain inherent in
  table-based ANS, with platform-specific BMI2 (amd64) and scalar (arm64)
  assembly kernels. We provide a formal latency model showing that the four-state
  design reduces amortised per-symbol latency from $L$ to $L/4$ cycles,
  achieving 66--142\% faster isolated FSE decompression with zero compression
  ratio penalty.

  \item \textbf{The simplicity thesis}: Empirical evidence that for images where
  a simple spatial predictor achieves $>$90\% of optimal decorrelation,
  additional transform complexity yields diminishing returns on lossless ratio
  while incurring proportional throughput cost. A $\sim$1,100-line
  Delta+RLE+FSE pipeline matches or beats JPEG~2000's wavelet+EBCOT on 7/8
  medical modalities---at dramatically higher throughput and a fraction of
  HTJ2K's implementation complexity.

  \item \textbf{Browser-native decoding for ubiquitous distribution}: A pure
  JavaScript ES module ($\sim$20\,KB, zero dependencies) and a Go WebAssembly
  build that decompress MIC files in any modern browser. PICS (Parallel Image
  Compressed Strips) extends this to multi-core browser decoding via
  \texttt{worker\_threads}, achieving up to 483\,MB/s on a 12-core workstation.
\end{enumerate}

We evaluate MIC on 21 clinical DICOM datasets spanning MR, CT, CR, X-ray,
mammography, nuclear medicine, radiography, secondary capture, and fluoroscopy
(XA) using fair in-process benchmarking via CGO bindings against HTJ2K
(OpenJPH), JPEG-LS (CharLS), and Delta+Zstandard. MIC achieves compression
ratios of 1.7$\times$--8.9$\times$. MIC-4state-C exceeds HTJ2K decompression
speed on 19/21 images single-threaded (ARM64) and 16/21 (AMD64, AVX2). PICS-C-8
exceeds HTJ2K on \textbf{all 21 images} on both platforms.
\end{abstract}

% ============================================================
\section{Introduction: The 16-Bit Alphabet Gap}
\label{sec:intro}

Medical imaging generates enormous volumes of data. A single digital breast
tomosynthesis (DBT) acquisition view produces 60 or more reconstructed
slices---e.g., 69 frames at $2457 \times 1890$ pixels, 10-bit depth on the
Hologic Selenia Dimensions---totalling over 600\,MB of raw pixel data per view.
A complete bilateral four-view screening study exceeds 2\,GB uncompressed.
Efficient lossless compression is essential for archival storage, network
transmission, and real-time rendering in clinical PACS and diagnostic viewers.

The DICOM standard~\cite{dicom} supports several transfer syntaxes for lossless
compression, including JPEG~2000~\cite{taubman2002}, JPEG-LS~\cite{loco1}, and
RLE Lossless. Each has well-documented tradeoffs between compression ratio and
decompression speed~\cite{liu2017,clunie2000}. But all share a less-discussed
limitation: \textbf{they were designed in an era of 8-bit data}.

JPEG's Huffman and arithmetic coders, JPEG-LS's Golomb-Rice coder, JPEG~2000's
MQ coder, and Zstandard's FSE backend all treat each coding unit as a small
fixed-size symbol (typically 4--8 bits). FSE as used in Zstandard~\cite{rfc8878}
caps its alphabet at 4,096 symbols. Deflate, LZ4, Brotli, and even
JPEG~XL~\cite{jpegxl} operate on 8-bit byte streams.

Medical images break this assumption. DICOM stores pixels at 10, 12, or 16 bits
per sample, yielding alphabet sizes of 1,024, 4,096, or 65,536 symbols. Even
after spatial delta prediction, the residual alphabet for a 16-bit CT image can
span tens of thousands of distinct values. An entropy coder that assumes an
8-bit alphabet must either:
\begin{itemize}
  \item \textbf{Re-encode 16-bit symbols as byte pairs}, losing inter-byte
  correlation and doubling the number of coding steps, or
  \item \textbf{Truncate the alphabet} at 256 and fall back to raw storage for
  rare symbols.
\end{itemize}

This mismatch has measurable consequences. We show (Section~\ref{sec:eval}) that
applying Zstandard to delta-encoded 16-bit medical images loses 10--22\% of
achievable compression compared to a 16-bit-native approach. The loss is
structural. A run of 1,000 identical 16-bit zero residuals is a single same-run
record in 16-bit RLE, but 2,000 bytes with no obvious structure to a byte-level
LZ77 matcher.

\subsection*{Contributions}

\begin{enumerate}
  \item We expose and quantify the 16-bit alphabet gap and demonstrate a
  16-bit-native RLE+FSE pipeline that closes it (Section~\ref{sec:pipeline}).
  \item We formalise multi-state ANS decoding as an ILP design axis, providing a
  latency model and platform-specific assembly implementations that achieve
  66--142\% speedup with no ratio penalty (Section~\ref{sec:multistate}).
  \item We present the simplicity thesis---empirical evidence that a simple
  spatial predictor with 16-bit entropy coding outperforms or matches
  wavelet+context-adaptive approaches at dramatically higher throughput
  (Section~\ref{sec:eval}).
  \item We demonstrate browser-native decoding as a distribution primitive: a
  $\sim$20\,KB JavaScript decoder and PICS achieve 483\,MB/s in a browser
  (Section~\ref{sec:browser}).
\end{enumerate}

% ============================================================
\section{Background and Related Work}
\label{sec:background}

\subsection{Lossless Image Compression in DICOM}

\textbf{JPEG~2000}~\cite{taubman2002} employs a 5/3 reversible integer wavelet
transform with embedded block coding (EBCOT). It achieves excellent compression
ratios but imposes significant decompression overhead. \textbf{High-Throughput
JPEG~2000 (HTJ2K)}~\cite{htj2k2019} replaces EBCOT with a faster block coder,
substantially improving decode speed.

\textbf{JPEG-LS}~\cite{loco1} uses context-adaptive prediction followed by
Golomb-Rice coding. It provides fast encoding with moderate compression ratios.

\textbf{RLE Lossless} (DICOM TS 1.2.840.10008.1.2.5) uses byte-level PackBits
run-length encoding. It is simple but typically achieves $<$1.5$\times$
compression because it operates on bytes rather than 16-bit pixel values.

\subsection{Entropy Coding and the Byte Assumption}

Asymmetric numeral systems (ANS)~\cite{duda2009,duda2013} generalise arithmetic
coding into a single-state machine with integer state transitions. Finite State
Entropy (FSE)~\cite{fse} is a table-driven tANS implementation that achieves
near-optimal compression with $\mathcal{O}(1)$ per-symbol operations. FSE was
popularised by Zstandard~\cite{rfc8878}. Recent theoretical analyses have
formalised ANS efficiency bounds~\cite{kosolobov2022} and statistical
interpretations~\cite{bamler2022}.

\textbf{All of these coders share one assumption: small alphabets.} Huffman
coding becomes impractical above $\sim$4,096 symbols. FSE in Zstandard caps at
4,096. The entire ecosystem was built for bytes.

\subsection{The Information-Theoretic Cost of Byte-Splitting}

Consider a 16-bit symbol $X$ decomposed into its high byte $X_H$ and low byte
$X_L$. The joint entropy satisfies:
\begin{equation}
  H(X) = H(X_H, X_L) = H(X_H) + H(X_L \mid X_H).
\end{equation}

A byte-level coder that processes $X_H$ and $X_L$ independently achieves at
best:
\begin{equation}
  H(X_H) + H(X_L) = H(X) + I(X_H;\, X_L),
\end{equation}
where $I(X_H;\, X_L)$ is the mutual information between the two bytes. This
mutual information is \emph{discarded}---it represents correlation that the
byte-level coder cannot exploit.

For medical image residuals after delta prediction, this mutual information is
substantial. When a residual is small (say, $\pm 30$), the high byte is always
\texttt{0x00} and carries essentially no information. On our test images,
$I(X_H;\, X_L)$ of delta residuals ranges from 0.3\,bits/symbol (MR) to
1.1\,bits/symbol (MG3), corresponding to 8--22\% of the total entropy---closely
matching the 10--22\% empirical advantage of MIC over Delta+Zstandard.

\subsection{Color Transforms for RGB Modalities}

The YCoCg-R transform~\cite{ycocgr} is a reversible integer approximation of
RGB-to-YCbCr. It decorrelates RGB channels into luminance (Y) and chrominance
(Co, Cg) with no loss of precision, and is used in MIC for ultrasound and
visible-light images.

% ============================================================
\section{Pipeline Design: Delta + 16-Bit RLE + Extended FSE}
\label{sec:pipeline}

MIC processes 16-bit greyscale pixels through three sequential stages. Each
stage is simple enough to decode in a single sequential pass with minimal
branching.

\begin{figure}[htbp]
\centering
\begin{lstlisting}
Raw 16-bit Pixels
    v  pixel - avg(top, left)
Delta Encoding
    v  same / diff runs
Run-Length Encoding (16-bit native)
    v  table-driven ANS
FSE (tANS) Entropy Coding
    v
Compressed Bitstream
\end{lstlisting}
\caption{MIC compression pipeline.}
\label{fig:pipeline}
\end{figure}

\subsection{Delta Encoding}

The delta encoder computes the prediction residual for each pixel as:
\begin{equation}
  r_{i,j} = x_{i,j} - \left\lfloor \frac{x_{i-1,j} + x_{i,j-1}}{2}
  \right\rfloor,
\end{equation}
where $x_{i,j}$ is the pixel at row $i$, column $j$. Boundary pixels use only
the available neighbour; the top-left corner pixel is stored raw.

For 16-bit images, residuals can span $[-65535, +65535]$. Rather than widening
to 32 bits, MIC uses an \textbf{overflow delimiter} scheme derived from the
image's effective bit depth $d = \lceil\log_2(\max(x)+1)\rceil$:
\begin{itemize}
  \item Residuals with $|r| < 2^{d-1}-1$ are stored directly as uint16 values.
  \item Larger residuals are preceded by a delimiter value $2^d - 1$, followed
  by the raw pixel value.
\end{itemize}

\textbf{Overflow delimiter: no collision by construction.} In-range residuals
are stored as \texttt{uint16(deltaThreshold + diff)}, mapping to the interval
$[1,\, 2^d-3]$. The delimiter is $D = 2^d - 1$. Since $2^d-3 < 2^d-1$, a
direct residual can never equal $D$; the two ranges are disjoint.

\textbf{Boundary-check-free interior loop.} The delta decompressor uses four
separate code paths for the corner pixel, first-row pixels, first-column pixels,
and interior pixels. The interior loop handles $(w-1)(h-1)$ out of $wh$ pixels
without any per-pixel boundary checks, leaving only the highly-predictable
overflow delimiter branch.

\textbf{Why not the MED predictor?} We implemented the JPEG-LS MED predictor
through the full RLE+FSE pipeline. Result: geometric mean improvement of only
$+0.9\%$ ratio, with MG4 actually regressing by $-1.7\%$, at 1.5--2$\times$
slower decompression due to the three-way conditional branch and diagonal
neighbour dependency.

\subsection{Run-Length Encoding (16-Bit Native)}

After delta coding, medical images produce long runs of identical residual
values (especially zero). MIC's RLE stage encodes the symbol stream as
alternating \textbf{same runs} and \textbf{diff runs}:
\begin{itemize}
  \item \textbf{Same run}: a count $c$ followed by a single value repeated
  $c$ times.
  \item \textbf{Diff run}: a count $c$ followed by $c$ distinct values.
\end{itemize}

\textbf{Why 16-bit RLE matters.} A run of 1,000 identical 16-bit zero residuals
is a single same-run record in MIC's RLE (2 uint16 words). The same data viewed
as 2,000 bytes has no obvious run structure to a byte-level compressor.

The run headers use a midpoint protocol: counts $\leq$ \texttt{midCount} signal
same runs; counts $>$ \texttt{midCount} signal diff runs.

\textbf{Same-run fast path.} Same-value runs---the most common pattern after
delta coding, often $>$80\% of all runs on smooth medical images---are
fast-pathed in the decoder to return the cached value without touching the
compressed-data slice. Each output symbol in a same run requires only a counter
decrement.

\subsection{FSE Entropy Coding (Extended to 65,535 Symbols)}

The final stage encodes the RLE output using tANS. \textbf{MIC extends FSE to
support up to 65,535 distinct symbols} (versus 4,095 in Zstandard), which is
necessary for 16-bit medical images where delta residuals can span a wide range.

\textbf{Adaptive table sizing.} MIC automatically adjusts the FSE table log via
thresholds derived from $\textit{symbolLen}$ and $\textit{symbolDensity} =
\textit{inputLen} / \textit{symbolLen}$:
\begin{itemize}
  \item Default: $\textit{tableLog} = 11$
  \item Bumped to 12 when $\textit{symbolDensity} > 32$ and
  $\textit{symbolLen} > 128$
  \item Bumped to 13 when $\textit{symbolLen} > 512$ and
  $\textit{symbolDensity} > 16$
\end{itemize}
This provides 4--7\% better compression on CR and mammography images
(Table~\ref{tab:tablelog}).

\begin{table}[htbp]
\centering
\footnotesize
\setlength{\tabcolsep}{4pt}
\caption{Effect of adaptive tableLog (11$\to$12).}
\label{tab:tablelog}
\begin{tabular}{lrrr}
\toprule
Image & tableLog=11 & Adaptive & Gain \\
\midrule
CR  & 3.47$\times$ & 3.63$\times$ & $+$4.4\% \\
MG1 & 7.99$\times$ & 8.57$\times$ & $+$7.1\% \\
MG2 & 7.98$\times$ & 8.55$\times$ & $+$7.1\% \\
\bottomrule
\end{tabular}
\end{table}

\textbf{Dynamic working-set sizing.} Encode and decode tables are sized to the
actual symbol range rather than the theoretical 65,536 maximum. For 8-bit images
this reduces the working set from 512\,KB to approximately 2\,KB, substantially
improving cache utilisation.

\subsection{Pipeline Selection Guidelines}

\begin{table}[htbp]
\centering
\footnotesize
\setlength{\tabcolsep}{4pt}
\caption{Recommended pipeline by use case.}
\label{tab:selection}
\begin{tabular}{p{3.5cm}p{3.8cm}}
\toprule
Condition & Pipeline \\
\midrule
Eff.\ bit depth $\leq 12$, speed & Delta+RLE+FSE (4-state) \\
Eff.\ bit depth $\leq 12$, ratio & Wavelet V2 SIMD \\
Full 16-bit range (CT) & Delta+RLE+FSE (avoid wavelet) \\
$\geq 0.5$\,MB, multi-core avail. & Add PICS (2--8 strips) \\
Multi-frame sequence & MIC2 independent mode \\
Single-frame RGB (US, VL) & MICR (YCoCg-R, no tiling) \\
Lossy/progressive & Use HTJ2K (MIC is lossless only) \\
\bottomrule
\end{tabular}
\end{table}

% ============================================================
\section{Multi-State ANS: ILP as a First-Class Design Axis}
\label{sec:multistate}

\subsection{The Serial Dependency Problem}

The standard FSE decode loop has a serial carried-dependency chain:
\begin{align}
  s[n{+}1] &= \texttt{decTable}[s[n]].\textit{newState} \nonumber\\
            &\quad + \texttt{getBits}\!\left(\texttt{decTable}[s[n]].\textit{nbBits}\right).
\end{align}
With a table-lookup latency of approximately $L \approx 4$ cycles (L1 cache
hit), the loop is limited to roughly one symbol per 4 cycles, regardless of
available execution units. Modern CPUs have 4--6 execution ports; the
single-state decoder uses exactly \textbf{one} dependency chain, leaving
75--83\% of the hardware capacity idle.

\subsection{Multi-State Interleaved Decoding}

MIC runs four independent state machines on interleaved symbol positions:
\begin{align*}
  \text{symbol}[0] &\leftarrow s_A, \quad \text{symbol}[4] \leftarrow s_A, \;\ldots \\
  \text{symbol}[1] &\leftarrow s_B, \quad \text{symbol}[5] \leftarrow s_B, \;\ldots \\
  \text{symbol}[2] &\leftarrow s_C, \quad \text{symbol}[6] \leftarrow s_C, \;\ldots \\
  \text{symbol}[3] &\leftarrow s_D, \quad \text{symbol}[7] \leftarrow s_D, \;\ldots
\end{align*}

\textbf{Correctness.} The encoder writes symbols at positions $0,4,8,\ldots$
into state $A$'s bitstream; positions $1,5,9,\ldots$ into state $B$'s; etc. The
decoder reverses this assignment. Each chain operates on a strict subset of
positions with no data dependency on other chains; correctness follows from the
correctness of single-state FSE applied independently to each subset.

\textbf{Latency model.}

\begin{table}[htbp]
\centering
\footnotesize
\setlength{\tabcolsep}{4pt}
\caption{Amortised latency and speedup for $k$-state FSE decoder.}
\label{tab:latency}
\begin{tabular}{clrr}
\toprule
States & Amortised latency & Theoretical & Empirical \\
\midrule
1 & $L$\;cy/sym   & 1.0$\times$ & 1.0$\times$ \\
2 & $L/2$\;cy/sym & 2.0$\times$ & 1.3$\times$ \\
4 & $L/4$\;cy/sym & 4.0$\times$ & 2.0$\times$ \\
\bottomrule
\end{tabular}
\end{table}

\textbf{Stream format.} The compressed stream is prefixed with a magic sequence
(\texttt{0xFF}, \texttt{0x02} for two-state; \texttt{0xFF}, \texttt{0x04} for
four-state) followed by a 4-byte symbol count, enabling the auto-detect
dispatcher to route streams transparently.

\subsection{Platform-Specific Assembly Kernels}

\textbf{AMD64 (BMI2).} Uses \texttt{SHLXQ}/\texttt{SHRXQ} for 3-operand
variable shifts, enabling four independent bit-extraction sequences per
iteration without CL register contention.

\textbf{ARM64.} Uses \texttt{LSLV}/\texttt{LSRV} variable-shift instructions
which natively take the count from any general-purpose register.

\textbf{AMD64 AVX2 --- wavelet kernels.} \texttt{wt53PredictAVX2} and
\texttt{wt53UpdateAVX2} process 8 \texttt{int32} values per cycle using
\texttt{VPADDD}/\texttt{VPSUBD}/\texttt{VPSRAD}, operating on contiguous arrays
produced by the blocked column layout (8 columns per cache line).

\textbf{Fallback.} Pure-Go implementations of all routines provide identical
functionality on non-amd64/arm64 targets (WebAssembly, RISC-V, etc.).

\subsection{Isolated FSE Decompression Throughput}

\begin{table}[htbp]
\centering
\footnotesize
\setlength{\tabcolsep}{4pt}
\caption{Isolated FSE decompression throughput (MB/s) for 1-, 2-, and 4-state
decoders (Intel Xeon @ 2.80\,GHz). MB/s over uncompressed RLE symbol stream.}
\label{tab:fse-throughput}
\begin{tabular}{lrrrr}
\toprule
Image & 1-state & 2-state & 4-state & 4 vs.\ 1 \\
\midrule
MR  & 164 & 207 & 298  & $+$82\%  \\
CT  & 113 & 126 & 195  & $+$73\%  \\
CR  & 177 & 182 & 310  & $+$75\%  \\
XR  & 193 & 172 & 320  & $+$66\%  \\
MG1 & 158 & 207 & 380  & $+$140\% \\
MG3 & 576 & 890 & 1,343 & $+$133\% \\
MG4 & 256 & 321 & 620  & $+$142\% \\
\bottomrule
\end{tabular}
\end{table}

% ============================================================
\section{Browser-Native Decoding}
\label{sec:browser}

MIC ships two browser decoder implementations:

\begin{table}[htbp]
\centering
\footnotesize
\setlength{\tabcolsep}{4pt}
\caption{Browser decoder options.}
\label{tab:browser-decoders}
\setlength{\tabcolsep}{2pt}
\begin{tabular}{p{2.8cm}rll}
\toprule
Decoder & Size & Deps. & Build \\
\midrule
\texttt{mic-decoder.js} & $\sim$20\,KB & None & None \\
\texttt{mic-decoder.wasm} & $\sim$2.5\,MB & \texttt{wasm\_exec.js} & \texttt{go build} \\
\bottomrule
\end{tabular}
\end{table}

The JavaScript decoder is a complete, self-contained ES module with zero npm
dependencies, working in any browser since 2020 (Chrome 67+, Firefox 68+,
Safari 14+).

\textbf{Single-threaded throughput} (Apple M2 Max, Node.js v24.8):

\begin{table}[htbp]
\centering
\footnotesize
\setlength{\tabcolsep}{4pt}
\caption{JavaScript decoder throughput (4-state), Node.js v24.8, M2 Max.}
\label{tab:js-perf}
\begin{tabular}{lrrrr}
\toprule
Image & Dimensions & Ratio & ms & MB/s \\
\midrule
MR  & $256{\times}256$   & 2.35$\times$ &   3.0 &  42 \\
CT  & $512{\times}512$   & 2.24$\times$ &  13.5 &  37 \\
CR  & $1760{\times}2140$ & 3.69$\times$ & 161.2 &  45 \\
MG1 & $1996{\times}2457$ & 8.79$\times$ &  80.9 & 116 \\
MG3 & $3064{\times}4774$ & 2.29$\times$ & 638.4 &  44 \\
\bottomrule
\end{tabular}
\end{table}

\textbf{Parallel decoding via PICS + \texttt{worker\_threads}.} PICS files
encode independent strips; each worker receives its strip, decompresses it
independently, and returns the pixel buffer.

\begin{table}[htbp]
\centering
\footnotesize
\setlength{\tabcolsep}{4pt}
\caption{PICS parallel JS decoder (\texttt{worker\_threads}), Apple M2 Max.}
\label{tab:pics-js}
\begin{tabular}{lrrrr}
\toprule
Image & Strips & ms & MB/s & Speedup \\
\midrule
MR ($256^2$)       & 4 &  1.3 &  94 & 2.57$\times$ \\
CT ($512^2$)       & 4 &  5.0 & 101 & 2.95$\times$ \\
CR ($1760{\times}2140$) & 8 & 31.7 & 227 & 5.53$\times$ \\
MG1 ($1996{\times}2457$) & 8 & 19.4 & 483 & 4.36$\times$ \\
\bottomrule
\end{tabular}
\end{table}

MG1 (a 9.35\,MB mammography image) decompresses in \textbf{19.4\,ms} at
483\,MB/s using 12 browser workers---well into real-time territory for
diagnostic viewing.

\textbf{Distribution model.} A web-based DICOM viewer can fetch a \texttt{.mic}
file directly from object storage (S3, GCS) and decode it client-side---zero
server-side compute, zero transcoding latency, works offline, works in service
workers. This is not possible today with HTJ2K or JPEG-LS.

% ============================================================
\section{Evaluation}
\label{sec:eval}

\subsection{Benchmarking Methodology and Dataset}

\textbf{Methodology.} All comparisons use in-process CGO library calls for both
MIC and competing codecs, measured on the same hardware in the same process
(\texttt{BenchmarkAllCodecs}, \texttt{-benchtime=10x}).

\textbf{Platforms.} Apple M2 Max (12-core ARM64) and Intel Core Ultra 9 285K
(24 P-core AMD64, \texttt{-march=native}).

\textbf{Dataset.} 21 clinical DICOM images spanning 10 modalities, drawn from
publicly available de-identified datasets (NEMA WG-04 compsamples, NEMA 1997
CD, Clunie DBT Case 1). See Table~\ref{tab:dataset}.

\begin{table}[htbp]
\centering
\footnotesize
\setlength{\tabcolsep}{4pt}
\caption{Test dataset --- 21 clinical DICOM images spanning 10 modalities.}
\label{tab:dataset}
\begin{tabular}{llrr}
\toprule
Image & Modality & Dimensions & Raw (MB) \\
\midrule
MR   & Brain MRI & $256{\times}256$ & 0.13 \\
CT   & CT scan   & $512{\times}512$ & 0.50 \\
CR   & Comp.\ radiography & $2140{\times}1760$ & 7.18 \\
XR   & X-ray & $2048{\times}2577$ & 10.1 \\
MG1  & Mammography & $2457{\times}1996$ & 9.35 \\
MG2  & Mammography & $2457{\times}1996$ & 9.35 \\
MG3  & Mammography & $4774{\times}3064$ & 27.3 \\
MG4  & Mammography & $4096{\times}3328$ & 26.0 \\
CT1  & CT scan & $512{\times}512$ & 0.50 \\
CT2  & CT scan & $512{\times}512$ & 0.50 \\
MG-N & Mammography & $3064{\times}4664$ & 27.3 \\
MR1  & Brain MRI & $512{\times}512$ & 0.50 \\
MR2  & Brain MRI & $1024{\times}1024$ & 2.00 \\
MR3  & Brain MRI & $512{\times}512$ & 0.50 \\
MR4  & Brain MRI & $512{\times}512$ & 0.50 \\
NM1  & Nuclear medicine & $256{\times}1024$ & 0.50 \\
RG1  & Radiography & $1841{\times}1955$ & 6.86 \\
RG2  & Radiography & $1760{\times}2140$ & 7.18 \\
RG3  & Radiography & $1760{\times}1760$ & 5.91 \\
SC1  & Sec.~capture     & $2048{\times}2487$ & 9.71 \\
XA1  & Fluoroscopy (XA) & $1024{\times}1024$ & 2.00 \\
\bottomrule
\end{tabular}
\end{table}

\subsection{Compression Ratios}

\begin{table*}[t]
\centering
\footnotesize
\setlength{\tabcolsep}{4pt}
\caption{Lossless compression ratios --- all codec variants, 21 images. Bold = best per row among non-JPEG-LS codecs.}
\label{tab:ratios}
\begin{tabular}{lrrrrrrr}
\toprule
Image & Raw (MB) & MIC & Wavelet & PICS-4 & PICS-8 & HTJ2K & JPEG-LS \\
\midrule
MR   & 0.13 & 2.35$\times$ & 2.38$\times$ & 2.28$\times$ & 2.21$\times$ & 2.38$\times$ & \textbf{2.52$\times$} \\
CT   & 0.50 & \textbf{2.24$\times$} & 1.67$\times$ & 2.15$\times$ & 1.96$\times$ & 1.77$\times$ & 2.68$\times$ \\
CR   & 7.18 & 3.69$\times$ & 3.81$\times$ & 3.70$\times$ & 3.71$\times$ & 3.77$\times$ & \textbf{3.96$\times$} \\
XR   & 10.1 & 1.74$\times$ & \textbf{1.76$\times$} & 1.75$\times$ & \textbf{1.76$\times$} & 1.67$\times$ & \textbf{1.76$\times$} \\
MG1  & 9.35 & 8.79$\times$ & 8.67$\times$ & 8.84$\times$ & \textbf{8.87$\times$} & 8.25$\times$ & 8.91$\times$ \\
MG2  & 9.35 & 8.77$\times$ & 8.65$\times$ & 8.83$\times$ & \textbf{8.85$\times$} & 8.24$\times$ & 8.90$\times$ \\
MG3  & 27.3 & 2.24$\times$ & 2.32$\times$ & 2.31$\times$ & \textbf{2.34$\times$} & 2.22$\times$ & 2.38$\times$ \\
MG4  & 26.0 & 3.47$\times$ & 3.59$\times$ & 3.59$\times$ & \textbf{3.62$\times$} & 3.51$\times$ & 3.71$\times$ \\
CT1  & 0.50 & \textbf{2.79$\times$} & 2.49$\times$ & 2.54$\times$ & 2.29$\times$ & 2.70$\times$ & 3.19$\times$ \\
CT2  & 0.50 & \textbf{3.49$\times$} & 2.87$\times$ & 3.11$\times$ & 2.72$\times$ & 3.29$\times$ & 4.54$\times$ \\
MG-N & 27.3 & 2.24$\times$ & 2.32$\times$ & 2.31$\times$ & \textbf{2.34$\times$} & 2.23$\times$ & 2.38$\times$ \\
MR1  & 0.50 & 2.09$\times$ & 2.14$\times$ & 2.10$\times$ & 2.08$\times$ & 2.13$\times$ & \textbf{2.30$\times$} \\
MR2  & 2.00 & 3.28$\times$ & 3.34$\times$ & 3.31$\times$ & 3.31$\times$ & 3.35$\times$ & \textbf{3.52$\times$} \\
MR3  & 0.50 & 3.93$\times$ & 4.09$\times$ & 3.89$\times$ & 3.84$\times$ & 4.33$\times$ & \textbf{4.51$\times$} \\
MR4  & 0.50 & 4.12$\times$ & 4.18$\times$ & 4.09$\times$ & 4.03$\times$ & 4.21$\times$ & \textbf{4.49$\times$} \\
NM1  & 0.50 & 5.15$\times$ & 5.02$\times$ & 5.26$\times$ & \textbf{5.28$\times$} & 5.76$\times$ & 6.28$\times$ \\
RG1  & 6.86 & 1.70$\times$ & 1.70$\times$ & 1.70$\times$ & 1.69$\times$ & 1.63$\times$ & \textbf{1.72$\times$} \\
RG2  & 7.18 & 4.23$\times$ & 4.32$\times$ & 4.28$\times$ & 4.30$\times$ & 4.32$\times$ & \textbf{4.51$\times$} \\
RG3  & 5.91 & 6.08$\times$ & 6.82$\times$ & 6.11$\times$ & 6.12$\times$ & 6.99$\times$ & \textbf{7.31$\times$} \\
SC1  & 9.71 & 3.71$\times$ & 3.70$\times$ & 3.73$\times$ & 3.74$\times$ & 3.85$\times$ & \textbf{4.73$\times$} \\
XA1  & 2.00 & 5.01$\times$ & 4.94$\times$ & 5.04$\times$ & \textbf{5.03$\times$} & 4.88$\times$ & 5.39$\times$ \\
\bottomrule
\end{tabular}
\end{table*}

\textbf{Analysis.} Mammography achieves the highest ratios (up to 8.87$\times$)
due to smooth tissue regions producing near-zero RLE runs. CT with full 16-bit
dynamic range achieves 2.24$\times$ (MIC) vs.\ 1.67$\times$ (Wavelet): coefficient
overflow in the 5/3 lifting step inflates compressed size. PICS improves ratio on
large images because each strip receives a dedicated FSE table that specialises to
its local residual distribution:
$\sum_k |S_k| \cdot H(S_k) < |S| \cdot H(S)$.

\textbf{16-bit gap vs.\ Zstandard.} MIC outperforms Delta+Zstandard (at level
19) by 10--22\% on all tested modalities. The gap is structural: mutual
information $I(X_H;\, X_L)$ of delta residuals ranges from 0.3\,bits/symbol
(MR) to 1.1\,bits/symbol (MG3), matching the empirical gap exactly.

\subsection{Decompression Speed --- ARM64}

MIC-4state-C is the recommended production variant. It is the fastest
single-threaded decoder on \textbf{19 of 21 images on ARM64}. PICS-C-8 exceeds
HTJ2K on \textbf{all 21 images}.

\begin{table*}[t]
\centering
\footnotesize
\setlength{\tabcolsep}{4pt}
\caption{Decompression throughput (MB/s), Apple M2 Max (ARM64), \texttt{-O3}. $\dagger$ = image too small for PICS-C-8; use PICS-C-4.}
\label{tab:decomp-arm}
\begin{tabular}{lrrrrrrrr}
\toprule
Image & MIC-Go & MIC-4s-C & Wvlt+SIMD & HTJ2K & JPEG-LS & PICS-C-2 & PICS-C-4 & PICS-C-8 \\
\midrule
MR   & 144 & \textbf{348} & 248 & 265 & 102 & 530 & \textbf{710}$^\dagger$ & 482 \\
CT   & 191 & \textbf{356} & 316 & 307 & 137 & 524 & 955 & \textbf{1,092} \\
CR   & 296 & 524 & \textbf{567} & 367 & 153 & 867 & 1,635 & \textbf{2,661} \\
XR   & 308 & 533 & \textbf{627} & 334 & 108 & 874 & 1,666 & \textbf{3,025} \\
MG1  & 482 & 683 & 678 & \textbf{810} & 409 & 1,205 & 2,112 & \textbf{3,656} \\
MG2  & 479 & 686 & 697 & \textbf{790} & 416 & 1,225 & 2,120 & \textbf{3,773} \\
MG3  & 308 & \textbf{531} & 422 & 338 & 153 & 864 & 1,673 & \textbf{3,117} \\
MG4  & 417 & \textbf{625} & 516 & 548 & 184 & 1,093 & 2,004 & \textbf{3,689} \\
CT1  & 239 & \textbf{436} & 425 & 362 & 182 & 686 & 1,013 & \textbf{1,183} \\
CT2  & 238 & 439 & \textbf{481} & 375 & 175 & 676 & 1,041 & \textbf{1,189} \\
MG-N & 316 & \textbf{536} & 468 & 340 & 153 & 883 & 1,711 & \textbf{3,175} \\
MR1  & 278 & \textbf{521} & 435 & 325 & 116 & 751 & 1,207 & \textbf{1,402} \\
MR2  & 333 & \textbf{563} & 498 & 388 & 172 & 913 & 1,552 & \textbf{2,466} \\
MR3  & 375 & \textbf{639} & 507 & 441 & 236 & 908 & 1,430 & \textbf{1,614} \\
MR4  & 316 & \textbf{571} & 479 & 406 & 197 & 818 & 1,341 & \textbf{1,558} \\
NM1  & 327 & \textbf{632} & 575 & 410 & 210 & 888 & 1,400 & \textbf{1,679} \\
RG1  & 235 & 406 & \textbf{584} & 332 & 104 & 602 & 1,128 & \textbf{2,017} \\
RG2  & 367 & 590 & \textbf{644} & 443 & 193 & 986 & 1,803 & \textbf{3,194} \\
RG3  & 374 & 604 & \textbf{656} & 562 & 246 & 1,035 & 1,944 & \textbf{3,302} \\
SC1  & 375 & \textbf{587} & 388 & 401 & 229 & 1,017 & 1,861 & \textbf{3,279} \\
XA1  & 331 & \textbf{576} & 459 & 419 & 204 & 928 & 1,583 & \textbf{2,493} \\
\bottomrule
\end{tabular}
\end{table*}

MIC-4state-C leads 13/21 images single-threaded. \textbf{MIC-4state-C is
1.7--5.0$\times$ faster than JPEG-LS on all 21 images.} PICS-C-8 peaks at
3,773\,MB/s on MG2, enabling sub-millisecond decompression of large mammography
images on a 12-core workstation.

\subsection{Decompression Speed --- AMD64}

On AMD64 (\texttt{-O3 -march=native}), \texttt{MIC-4state-SIMD} activates the
BMI2/PDEP decoder, and Wavelet+SIMD gains AVX2 kernels.

\begin{table*}[t]
\centering
\footnotesize
\setlength{\tabcolsep}{4pt}
\caption{Decompression throughput (MB/s), Intel Core Ultra 9 285K (AMD64), \texttt{-O3 -march=native}. $\dagger$ = image too small for PICS-C-8.}
\label{tab:decomp-amd64}
\begin{tabular}{lrrrrrrr}
\toprule
Image & MIC-Go & MIC-4s-C & MIC-4s-SIMD & Wvlt+SIMD & HTJ2K & PICS-C-4 & PICS-C-8 \\
\midrule
MR   & 251 & 501 & \textbf{708} & 194 & 570 & 301 & 124$^\dagger$ \\
CT   & 231 & 403 & 487 & 364 & \textbf{544} & \textbf{676} & 339 \\
CR   & 324 & 599 & \textbf{744} & 723 & 708 & 1,738 & \textbf{2,435} \\
XR   & 363 & 601 & \textbf{803} & 681 & 570 & 1,714 & \textbf{1,994} \\
MG1  & 617 & 789 & 1,119 & 746 & \textbf{1,235} & 1,872 & \textbf{2,514} \\
MG2  & 562 & 800 & 1,166 & 848 & \textbf{1,297} & \textbf{2,244} & 2,085 \\
MG3  & 357 & 633 & 669 & \textbf{678} & 644 & 1,823 & \textbf{2,538} \\
MG4  & 523 & 743 & 773 & \textbf{808} & 916 & 2,124 & \textbf{2,707} \\
CT1  & 322 & 520 & \textbf{676} & 525 & 657 & \textbf{857} & 423 \\
CT2  & 293 & 514 & 636 & \textbf{705} & 627 & \textbf{847} & 687 \\
MG-N & 368 & 635 & 669 & \textbf{705} & 643 & 1,872 & \textbf{2,191} \\
MR1  & 356 & 609 & \textbf{766} & 532 & 654 & \textbf{945} & 366 \\
MR2  & 352 & 658 & \textbf{975} & 601 & 749 & \textbf{1,121} & 793 \\
MR3  & 450 & 728 & \textbf{809} & 676 & 802 & \textbf{920} & 705 \\
MR4  & 390 & 596 & \textbf{861} & 738 & 660 & 493 & \textbf{701} \\
NM1  & 367 & \textbf{717} & 668 & 715 & 627 & \textbf{783} & 405 \\
RG1  & 302 & 497 & 593 & \textbf{705} & 557 & 1,248 & \textbf{1,494} \\
RG2  & 433 & 676 & \textbf{861} & 811 & 823 & 1,841 & \textbf{1,912} \\
RG3  & 460 & 706 & 858 & 881 & \textbf{894} & \textbf{1,969} & 1,613 \\
SC1  & 464 & 695 & \textbf{888} & 710 & 728 & 1,819 & \textbf{1,889} \\
XA1  & 413 & 693 & \textbf{836} & 538 & 797 & 1,173 & \textbf{1,513} \\
\bottomrule
\end{tabular}
\end{table*}

MIC-4state-SIMD leads 16/21 images single-threaded; HTJ2K leads on MG1, MG2,
MG4, CT, and RG3. \textbf{PICS-C-8 exceeds HTJ2K on all 21 images} on both
platforms.

\subsection{Encoding Speed}

On ARM64, MIC-4state-C achieves 243--644\,MB/s single-threaded (1.2--2.5$\times$
faster than MIC-Go). PICS-8 scales to \textbf{1.2--2.1\,GB/s} on large images
using 8 goroutines, sufficient for real-time PACS ingestion. HTJ2K encode is
comparable to MIC-Go on most images; JPEG-LS is consistently 2--4$\times$ slower
to encode. Wavelet+SIMD encode is 2--4$\times$ slower than MIC-Go due to the
multi-level forward transform cost.

\subsection{Multi-Frame and Single-Frame RGB Compression}

\textbf{Multi-frame (MIC2)} --- 69-frame breast tomosynthesis ($2457 \times
1890 \times 69$, 614\,MB raw):

\begin{table}[htbp]
\centering
\footnotesize
\setlength{\tabcolsep}{4pt}
\caption{69-frame breast tomosynthesis compression.}
\label{tab:multiframe}
\begin{tabular}{lrr}
\toprule
Mode & Compressed & Ratio \\
\midrule
Independent (spatial) & 46.1\,MB & \textbf{13.3$\times$} \\
Temporal (inter-frame) & 47.5\,MB & 12.9$\times$ \\
\bottomrule
\end{tabular}
\end{table}

Independent mode outperforms temporal mode on this dataset because
tomosynthesis frames differ in X-ray projection angle, not temporal motion.
Temporal mode is designed for genuine temporal sequences (cardiac cine MRI,
fluoroscopy).

\textbf{Single-Frame RGB (MICR)} --- YCoCg-R + Delta+RLE+FSE applied to the
full image (NEMA compsamples):

\begin{table}[htbp]
\centering
\footnotesize
\setlength{\tabcolsep}{4pt}
\caption{Single-frame RGB compression ratios (MICR, NEMA compsamples).}
\label{tab:rgb}
\begin{tabular}{lrrp{2.0cm}}
\toprule
Image & Dimensions & Ratio & Notes \\
\midrule
US1    & $640{\times}480$   & 6.24$\times$ & Uniform bg \\
VL1--3 & $756{\times}486$   & 3.2--3.5$\times$ & Visible light \\
VL4    & $2226{\times}1868$ & 1.86$\times$ & High detail \\
VL5    & $2670{\times}3340$ & 1.56$\times$ & Skin texture \\
VL6    & $756{\times}486$   & 1.93$\times$ & --- \\
\bottomrule
\end{tabular}
\end{table}

% ============================================================
\section{Discussion}
\label{sec:discussion}

\subsection{Signal Processing and Lossless Compression}

For \textbf{lossy} compression, filter-bank decomposition is highly
effective---high-frequency subbands can be quantised with little perceptual
impact. For \textbf{lossless} compression the picture is more nuanced. The
evaluation results (Section~\ref{sec:eval})---particularly the CT wavelet
regression (1.67$\times$ vs.\ 2.24$\times$ for Delta)---are explained by three
structural mechanisms:

\begin{enumerate}
  \item \textbf{Coefficient overflow.} Even with integer lifting (the 5/3 Le
  Gall wavelet), the predict step can produce detail coefficients exceeding the
  input dynamic range. For 16-bit inputs, worst-case expansion per level is 3/2;
  for $L=5$ levels, $(3/2)^5 \approx 7.6\times$---well beyond uint16, forcing
  int32 promotion or escape coding.

  \item \textbf{The update step reintroduces correlation.} The lifting update
  adjusts low-pass coefficients to preserve the signal mean. For lossless coding
  of images already well-predicted by a simple linear filter, it adds residual
  energy to the low-pass band without a compensating reduction elsewhere.

  \item \textbf{Higher memory bandwidth.} The 2D wavelet inverse transform
  requires two passes, each visiting every element at int32 width (4\,bytes vs.\
  2\,bytes per sample).
\end{enumerate}

By contrast, simple spatial prediction leaves a single residual image whose
statistics are homogeneous across the frame---ideal for a uniform 16-bit entropy
coding scheme applied once.

\subsection{The Cost-of-Complexity Frontier}

\begin{table}[htbp]
\centering
\footnotesize
\setlength{\tabcolsep}{4pt}
\caption{Pareto frontier: ratio vs.\ decompression throughput (approx.\ geo.\ means).}
\label{tab:pareto}
\setlength{\tabcolsep}{2pt}
\begin{tabular}{lrrrp{1.5cm}}
\toprule
Codec & Ratio & Decomp. & LOC & Browser \\
\midrule
$\Delta$+RLE+FSE (Go)   & 3.12$\times$ & 310\,MB/s & 1.1k & \textbf{Yes (20\,KB)} \\
$\Delta$+RLE+FSE (4s-C) & 3.12$\times$ & 530\,MB/s & 1.5k & \textbf{Yes} \\
Wavelet V2 SIMD         & 3.28$\times$ & 780\,MB/s & 2.0k & Possible \\
HTJ2K (OpenJPH)         & 3.15$\times$ & 460\,MB/s & 20k  & No \\
JPEG-LS (CharLS)        & 3.44$\times$ & 130\,MB/s & 5.0k & No \\
$\Delta$+zstd-19        & 2.72$\times$ & 300\,MB/s & N/A  & Partial \\
\bottomrule
\end{tabular}
\end{table}

JPEG-LS trades 58\% of decompression speed for 10\% more compression---unfavorable
for clinical systems that decompress 10$\times$ more often than they compress.
HTJ2K's $\sim$20,000-line implementation is $\sim$18$\times$ larger than MIC's
Delta+RLE+FSE pipeline and has no practical browser decoder, making it unsuitable
for client-side web deployment.

\subsection{Limitations}

\begin{itemize}
  \item \textbf{No lossy compression}, progressive decoding, or
  region-of-interest coding for greyscale images. The wavelet pipeline provides
  a natural foundation for a future lossy mode.
  \item \textbf{Test dataset}: 21 images across 10 modalities from three public
  de-identified repositories. Large-scale benchmarks from public repositories
  (e.g., TCIA) would further strengthen generality claims.
  \item \textbf{Temporal prediction} is underperforming on the one available
  clinical dataset and needs evaluation on cardiac cine MRI and fluoroscopy.
  \item \textbf{Benchmark confidence intervals}: run-to-run variance is $<$2\%
  on M2 Max and $<$5\% on AWS instances; formal confidence intervals are deferred
  to future work.
\end{itemize}

\subsection{Pathway to Clinical Impact}

MIC could be registered as a DICOM Private Transfer Syntax, allowing PACS
vendors to adopt it without modifying the DICOM standard. The $\sim$20\,KB
JavaScript decoder enables a direct storage-and-serve distribution model:
compressed MIC files can be fetched from object storage and decoded client-side,
eliminating server-side transcoding.

% ============================================================
\section{Conclusion}
\label{sec:conclusion}

We have presented MIC, a lossless medical image codec built on a simple
observation: the compression ecosystem has a 30-year blind spot for 16-bit data.
By extending FSE to 65,535 symbols and pairing it with a 16-bit-native RLE
stage, MIC outperforms byte-oriented Zstandard by 10--22\% on all tested
modalities. The four-state interleaved ANS decoder demonstrates that entropy
coder design should target instruction-level parallelism width, not just coding
efficiency, achieving 66--142\% FSE speedup with no ratio penalty.

The simplicity thesis is validated empirically: a three-stage Delta+RLE+FSE
pipeline ($\sim$1,100 lines of Go) matches or beats HTJ2K on 18/21 medical
images for lossless compression ratio, at dramatically higher throughput.

Key results across 21 clinical DICOM images spanning 10 modalities:
\begin{itemize}
  \item \textbf{Compression ratios}: 1.7$\times$--8.9$\times$ (greyscale);
  1.56$\times$--6.24$\times$ (single-frame RGB US/VL).
  \item \textbf{vs.\ HTJ2K}: MIC-4state-C exceeds HTJ2K on 19/21 images
  single-threaded (ARM64); MIC-4state-SIMD exceeds on 16/21 (AMD64 AVX2);
  PICS-C-8 exceeds HTJ2K on \textbf{all 21 images} on both platforms.
  \item \textbf{vs.\ JPEG-LS}: 1.7--5.0$\times$ faster decompression at 1--30\%
  lower ratio depending on modality.
  \item \textbf{vs.\ Delta+Zstandard}: 10--22\% better compression on all
  modalities.
  \item \textbf{Browser decoder}: $\sim$20\,KB pure JS decoder; PICS +
  \texttt{worker\_threads} achieves 483\,MB/s (MG1) with no server-side
  transcoding.
  \item \textbf{Compact}: $\sim$1,100 lines of Go --- approximately $18\times$
  less than HTJ2K's $\sim$20,000-line implementation.
\end{itemize}

MIC is open-source: \url{https://github.com/pappuks/medical-image-codec}.

% ============================================================
\balance
\begin{thebibliography}{99}

\bibitem{dicom}
NEMA, ``Digital Imaging and Communications in Medicine (DICOM),''
\url{https://www.dicomstandard.org/}, 2024.

\bibitem{taubman2002}
D.\,S.\,Taubman and M.\,W.\,Marcellin,
\textit{JPEG2000: Image Compression Fundamentals, Standards and Practice}.
Springer, 2002.

\bibitem{htj2k2019}
D.\,S.\,Taubman, ``High throughput JPEG 2000 (HTJ2K): Algorithm, performance
evaluation, and potential,'' \textit{Proc.\ SPIE}, vol.\,11137, 2019.

\bibitem{loco1}
M.\,J.\,Weinberger, G.\,Seroussi, and G.\,Sapiro,
``The LOCO-I lossless image compression algorithm: Principles and
standardization into JPEG-LS,''
\textit{IEEE Trans.\ Image Process.}, vol.\,9, no.\,8, pp.\,1309--1324, 2000.

\bibitem{jpegxl}
J.\,Alakuijala \textit{et al.},
``JPEG XL next-generation image compression architecture and coding tools,''
\textit{Proc.\ SPIE 11137}, Sep.\,2019.

\bibitem{duda2009}
J.\,Duda, ``Asymmetric numeral systems,'' arXiv:0902.0271, 2009.

\bibitem{duda2013}
J.\,Duda, ``Asymmetric numeral systems: entropy coding combining speed of
Huffman coding with compression rate of arithmetic coding,''
arXiv:1311.2540, 2013.

\bibitem{fse}
Y.\,Collet, ``Finite State Entropy --- A new breed of entropy coder,''
\url{https://github.com/Cyan4973/FiniteStateEntropy}, 2013.

\bibitem{rfc8878}
Y.\,Collet and M.\,Kucherawy,
``Zstandard Compression and the `application/zstd' Media Type,''
RFC~8878, IETF, Feb.\,2021.

\bibitem{schwartz1964}
E.\,S.\,Schwartz and B.\,Kallick,
``Generating a canonical prefix encoding,''
\textit{Commun.\ ACM}, vol.\,7, no.\,3, pp.\,166--169, 1964.

\bibitem{mahapatra2007}
S.\,Mahapatra and K.\,Singh,
``An FPGA-based implementation of multi-alphabet arithmetic coding,''
\textit{IEEE Trans.\ Circuits Syst.\ I}, vol.\,54, no.\,8, pp.\,1678--1686, 2007.

\bibitem{kosolobov2022}
D.\,Kosolobov, ``Efficiency of ANS Entropy Encoders,''
arXiv:2201.02514, 2022.

\bibitem{bamler2022}
R.\,Bamler, ``Understanding Entropy Coding With Asymmetric Numeral Systems
(ANS): a Statistician's Perspective,'' arXiv:2201.01741, 2022.

\bibitem{giesen2014}
F.\,Giesen, ``Interleaved entropy coders,'' arXiv:1402.3392, 2014.

\bibitem{giesen2014b}
F.\,Giesen, ``ryg\_rans: Public domain rANS encoder/decoder,''
\url{https://github.com/rygorous/ryg_rans}, 2014.

\bibitem{lin2023}
F.\,Lin, K.\,Arunruangsirilert, H.\,Sun, and J.\,Katto,
``Recoil: Parallel rANS Decoding with Decoder-Adaptive Scalability,''
\textit{Proc.\ 52nd ICPP '23}, pp.\,31--40, 2023.

\bibitem{weissenberger2019}
A.\,Weissenberger and B.\,Schmidt,
``Massively Parallel ANS Decoding on GPUs,''
\textit{Proc.\ 48th ICPP '19}, Article~100, 2019.

\bibitem{wu1997}
X.\,Wu and N.\,Memon,
``Context-based, adaptive, lossless image coding,''
\textit{IEEE Trans.\ Commun.}, vol.\,45, no.\,4, pp.\,437--444, 1997.

\bibitem{sneyers2016}
J.\,Sneyers and P.\,Wuille,
``FLIF: Free Lossless Image Format based on MANIAC Compression,''
\textit{Proc.\ IEEE ICIP}, 2016.

\bibitem{legall1988}
D.\,Le Gall and A.\,Tabatabai,
``Sub-band coding of digital images using symmetric short kernel filters and
arithmetic coding techniques,'' \textit{Proc.\ IEEE ICASSP}, pp.\,761--764, 1988.

\bibitem{sweldens1996}
W.\,Sweldens,
``The Lifting Scheme: A custom-design construction of biorthogonal wavelets,''
\textit{Appl.\ Comput.\ Harmon.\ Anal.}, vol.\,3, no.\,2, pp.\,186--200, 1996.

\bibitem{daubechies1998}
I.\,Daubechies and W.\,Sweldens,
``Factoring wavelet transforms into lifting steps,''
\textit{J.\ Fourier Anal.\ Appl.}, vol.\,4, no.\,3, pp.\,247--269, 1998.

\bibitem{calderbank1998}
A.\,R.\,Calderbank, I.\,Daubechies, W.\,Sweldens, and B.-L.\,Yeo,
``Wavelet transforms that map integers to integers,''
\textit{Appl.\ Comput.\ Harmon.\ Anal.}, vol.\,5, no.\,3, pp.\,332--369, 1998.

\bibitem{ycocgr}
H.\,S.\,Malvar and G.\,J.\,Sullivan,
``YCoCg-R: A color space with RGB reversibility and low dynamic range,''
JVT Doc.\ JVT-I014, 2003.

\bibitem{liu2017}
F.\,Liu \textit{et al.},
``The Current Role of Image Compression Standards in Medical Imaging,''
\textit{Information}, vol.\,8, no.\,4, p.\,131, Oct.\,2017.

\bibitem{clunie2000}
D.\,A.\,Clunie,
``Lossless Compression of Grayscale Medical Images: Effectiveness of
Traditional and State-of-the-Art Approaches,''
\textit{Proc.\ SPIE Medical Imaging}, 2000.

\bibitem{taubman2022}
D.\,Taubman, A.\,Naman, M.\,Smith, P.-A.\,Lemieux, H.\,Saadat, O.\,Watanabe,
and R.\,Mathew,
``High Throughput JPEG~2000 for Video Content Production and Delivery Over IP
Networks,''
\textit{Front.\ Signal Process.}, vol.\,2, Article 885644, Apr.\,2022.

\bibitem{williams2009}
S.\,Williams, A.\,Waterman, and D.\,Patterson,
``Roofline: An insightful visual performance model for multicore
architectures,''
\textit{Commun.\ ACM}, vol.\,52, no.\,4, pp.\,65--76, Apr.\,2009.

\bibitem{openjph}
A.\,N.\,Aous,
``OpenJPH: An open-source implementation of High-Throughput JPEG~2000,''
\url{https://github.com/aous72/OpenJPH}, 2024.

\bibitem{charls}
Team CharLS,
``CharLS: A C/C++ JPEG-LS library implementation,''
\url{https://github.com/team-charls/charls}, 2024.

\bibitem{nohtj2k}
As of the submission date of this paper, no production-ready WebAssembly or
JavaScript decoder exists for High-Throughput JPEG~2000. The OpenJPH
repository~\cite{openjph} does not provide a WASM build or npm package.

\bibitem{fzstd}
N.\,Tindall \textit{et al.},
``fzstd: Pure JavaScript Zstandard decompressor,''
\url{https://github.com/101arrowz/fzstd}, 2024.

\end{thebibliography}

\end{document}
